\documentclass[main.tex]{subfiles}
\begin{document}
In the previous sections we proposed explicit combinatorial formulas for the last layer of the socle filtration of an indecomposable injective tensor $\mathfrak{sp}_\infty-$module, as well as for the last term of a minimal injective resolution of a simple tensor $\mathfrak{sp}_\infty-$module. Using the equivalence of the categories of tensor $\mathfrak{sp}_\infty-$ and $\mathfrak{so}_\infty-$modules, we applied our results also to tensor $\mathfrak{so}_\infty-$modules. \\

It makes sense to try to find explicit information about injective tensor $\mathfrak{sp}_\infty-$ and $\mathfrak{so}_\infty-$modules and injective resolutions. In particular, I. Penkov and V. Tsanov made the following conjecture: \\

Consider the socle layers of an injective tensor $\mathfrak{sp}_\infty-$module $\Gamma_\lambda,$ and let $s(k+1):= \sum\limits_{|\gamma|=k}\LR{\lambda}{\mu}{(2\gamma)^T},$ i.e., let $s(k+1)$ be the sum of the multiplicities of all simple constituents $\Gamma_{<\mu>}$ of $\soc{k+1}\Gamma_\lambda$. The conjecture claims that the socle filtration layers have a diamond shape structure, meaning that the inequalities
$$
    1=s(1)\leq s(2) \leq s(3) \leq ... \leq s(m) \geq s(m-1) \geq s(m-2) \geq ... \geq s(|\lambda|^{even}+1),
$$
hold for  $m = \lfloor{\frac{|\lambda|^{even}}{2}}\rfloor+1.$

\end{document}